% ============================================================
%  Differencing transformations: BSSD and DD
% ============================================================

PPP-RTK with the single-difference transform and the baseline modes do not consume
the undifferenced observations of \Cref{eq:obs-code,eq:obs-phase} directly.  Each
forms a linear combination across satellites (BSSD), receivers (BRSD), or both (DD).
The combination cancels common-mode terms and leaves a residual structure in which
the ambiguities can be integer.  This section gives the algebraic cancellations and
the matrix-form S-transforms that map undifferenced parameters to differenced ones.

\subsection{Between-satellite single difference (BSSD)}\label{sec:bssd}

Each constellation $g$ has a pivot satellite $\piv{g}$, elected by
\texttt{PivotManager::selectPivotCandidates} under a sticky retention policy (the
ordering is given in the \emph{Pivot selection} paragraph below).  The same
\texttt{PivotManager} pivot serves both the combined (ionosphere-free) and the
uncombined parameterisation.  For any non-pivot satellite $\sat \in \satset{g} \setminus \{\piv{g}\}$
and signal $\sigl$, define
\begin{align}
\bsd{\code^{\sat,\piv{g}}_{\rx,\sigl}}
&:= \code^\sat_{\rx,\sigl} - \code^{\piv{g}}_{\rx,\sigl}, \notag\\
\bsd{\phaseo^{\sat,\piv{g}}_{\rx,\sigl}}
&:= \phaseo^\sat_{\rx,\sigl} - \phaseo^{\piv{g}}_{\rx,\sigl}.
\end{align}
Substituting \Cref{eq:obs-code,eq:obs-phase} and cancelling the
receiver-side common terms ($\sLight\,\mathrm{d}t_\rx$,
$\sLight\,\delta t^{(g)}_\rx$, $B_{\rx,\sigl}$, $b_{\rx,\sigl}$) yields
\begin{empheq}[box=\fbox]{align}
\bsd{\code^{\sat,\piv{g}}_{\rx,\sigl}}
&\;=\; \bsd{\geomrange^{\sat,\piv{g}}_\rx}
     - \sLight\,\bsd{\mathrm{d}t^{\sat,\piv{g}}}
     + \bsd{T^{\sat,\piv{g}}_\rx}
     + \mu_\sigl\,\bsd{I^{\sat,\piv{g}}_{\rx,1}}
     - \bsd{B^{\sat,\piv{g}}_\sigl}
     + \bsd{\varepsilon^\code}, \label{eq:bssd-code}\\
\bsd{\phaseo^{\sat,\piv{g}}_{\rx,\sigl}}
&\;=\; \bsd{\geomrange^{\sat,\piv{g}}_\rx}
     - \sLight\,\bsd{\mathrm{d}t^{\sat,\piv{g}}}
     + \bsd{T^{\sat,\piv{g}}_\rx}
     - \mu_\sigl\,\bsd{I^{\sat,\piv{g}}_{\rx,1}}
     + \lambda_\sigl\,\bsd{N^{\sat,\piv{g}}_{\rx,\sigl}}
     - \bsd{b^{\sat,\piv{g}}_\sigl}
     + \bsd{\varepsilon^\phaseo}. \label{eq:bssd-phase}
\end{empheq}
The receiver-side cancellations are exact only when both observations enter the
filter with the same modelling (tropospheric mapping, antenna PCV correction and so
on).  The satellite-dependent terms do not cancel: $\bsd{T}$, for example, is the
difference of two slant delays at different elevations and is small but not zero.

\texttt{BssdTransform::apply} realises the cancellation by zeroing the design-matrix columns
of the receiver-shared states after differencing: receiver clock, phase clock, ISB,
clock rate and ISB rate (when allocated), and the receiver code-bias (DCB/IFB) states.
In both parameterisations it also zeroes the pivot's ambiguity columns (under the
NL/WL basis, the pivot's NL column and, on a non-reference band, its WL column).

The cancellation of $B_{\rx,\sigl}$ in \Cref{eq:bssd-code} holds for CDMA
constellations, whose group delay is a single per-signal receiver term.  GLONASS is
the exception.  Its FDMA inter-frequency biases are per-FCN-slot states, so two
satellites on different slots carry different receiver delays and a residual
$B^{(k(\sat))}_{\rx,\sigl} - B^{(k(\piv{g}))}_{\rx,\sigl}$ survives the single
difference.  \texttt{BssdTransform} therefore exempts the per-slot GLONASS IFB states
(identified by \texttt{StateKeys::isGlonassPerSlotIfbKey}) from the column zeroing.
\texttt{PivotManager::buildEpochMaps} drops GLONASS from the BSSD candidate maps
unless the flag is set, and \texttt{BssdTransform::apply} drops any row whose group
has no pivot, so with the flag off GLONASS contributes no single-differenced row.
GLONASS BSSD is off by default
(\texttt{estimation.between\_satellite\_sd.enable\_glonass}); when enabled,
\Cref{eq:bssd-code} acquires the surviving inter-slot IFB difference as an
additional estimated column.

\paragraph{Pivot selection.}
The incumbent pivot is kept while it remains at or above
\texttt{estimation.between\_satellite\_sd.pivot\_min\_elevation} (default
45 degrees), is not in the exclusion set, and passes the ephemeris-freshness bar
described below.  Otherwise a new pivot is elected.  The primary key is a hard
lexicographic tier on phase-signal coverage: a full-coverage candidate outranks any
partial-coverage candidate regardless of score, because the BSSD frame requires the
pivot to track every signal another satellite tracks.  Within a coverage tier the
additive score
$-\min\!\big(1,\sqrt{\sigma^2_{\hat N,\max}}\big) + \mathrm{el}/90 -
\tfrac12\,\mathbb{1}[\mathrm{el} < 20^\circ]$
ranks the candidates.  Float-ambiguity convergence $\sqrt{\sigma^2_{\hat N,\max}}$
and elevation enter as co-equal terms, with a low-elevation penalty below
$20^\circ$; a candidate with no ambiguity state yet takes the saturated convergence
penalty of 1.  The ephemeris-freshness bar excludes from pivot election any candidate
whose broadcast $|\text{obs} - \text{toe}|$ exceeds that of its constellation's
freshest qualified peer by more than \texttt{kPivotToeSpreadMaxS}
($\SI{1200}{s}$); the satellite remains usable as a non-pivot.  The bar prevents
the extrapolated orbit and clock error of a stale-ephemeris pivot from entering
every SD row it anchors.  Candidates below the minimum elevation are excluded; only
when no satellite of the constellation passes the filters does the election fall
back to the best-ranked or highest satellite without them.  The baseline DD path
scores differently: \texttt{DdPivotManager} keeps the incumbent while it remains in
the candidate set of satellites shared by both receivers, and otherwise ranks on
$\text{avg\_snr} + 0.01\cdot\text{num\_sigs}$, without elevation.

\paragraph{Pivot transitions.}
A handover has three kinds of trigger.  The first is selection: when the incumbent
fails the retention test above (it leaves the visible set, drops below the minimum
elevation, is excluded, or fails the freshness bar), $\piv{g}$ is re-elected from
the remaining satellites.  The second is a disturbed state:
\texttt{isPivotDisturbed} reports the pivot disturbed when any of its phase
ambiguity states has left $(0,0)$ (absolute value above \num{1e-6} or variance
above \num{1e-12}), for example after a cycle-slip reset.  The third is an event: the
Kalman filter's post-fit residual analyser calls
\texttt{PivotManager::requestForcedHandover} when the pivot's observation has
slipped while its filter state is intact, a case \texttt{isPivotDisturbed} cannot
detect.  The forced request carries an exclusion set, so the handover cannot land on
another offset satellite.  The \texttt{retired\_pivots\_this\_epoch\_} set bars a
pivot retired in the current epoch from re-election within it, which terminates
handover cascades, and the same-epoch re-entry gate
(\texttt{hasPendingForcedHandover}) applies the handover before the rebuilt
measurement update.  The request also carries a re-inflation set: non-pivot states
that absorbed a coherent pivot-side code bias have their covariance re-inflated
after the re-reference (when slip-state re-inflation is enabled), because the
information-conserving transform below cannot remove that offset.  The baseline DD
path has the counterpart \texttt{DdPivotManager::requestForcedRerank}, which drops
the incumbent and excludes the flagged satellites from the next re-rank.  The event
response is required because the pivot's states are frozen at zero variance (below):
a stiff, time-linked zero-$q$ quantity requires a wired jump response.

On a handover the float ambiguities
$\{\bsd{N^{\sat,\piv{g}^{\mathrm{old}}}_{\rx,\sigl}}\}_{\sat,\sigl}$
are re-expressed relative to $\piv{g}^{\mathrm{new}}$ by the linear
S-transformation $\stateV^{\mathrm{new}} = \bm{S}\,\stateV^{\mathrm{old}}$, with
$\bm{S}$ chosen so that
$\bsd{N^{\piv{g}^{\mathrm{new}},\piv{g}^{\mathrm{old}}}_{\rx,\sigl}}$
becomes the new reference (zero).  The covariance transforms as
$\covM^{\mathrm{new}} = \bm{S}\,\covM^{\mathrm{old}}\,\bm{S}^{\Tr}$.  The transform
conserves the filter's information and so preserves the integer information
accumulated in the ambiguities.  In the uncombined PPP-RTK parameterisation
\texttt{PivotManager} applies the same $\bm{S}\,\covM\,\bm{S}^{\Tr}$ re-reference to
three state groups (\texttt{reReferenceAmbiguities}, \texttt{reReferenceIonosphere}
and \texttt{reReferenceSatBiases}).  This is why \texttt{BssdTransform} zeroes the
pivot's ionosphere and sat-bias $\bm{H}$-columns after differencing, and why
$\bsd{I^{\sat,\piv{g}}_{\rx,1}}$ in \Cref{eq:bssd-code,eq:bssd-phase} is a
pivot-referenced SD quantity rather than an absolute slant delay.

Where the transform cannot carry a state onto the new pivot, a re-seed follows or
replaces it.  In PPP-RTK three cases arise.  An ambiguity on a phase signal the new pivot does not
track stays referenced to the old pivot, so it is reset and marked for re-seeding
against the new pivot.  When the old pivot was disturbed, its re-referenced
ambiguities are re-seeded from the single-differenced observations.  When a
disturbed pivot has no eligible replacement (every other satellite of the
constellation is excluded or absent),
\texttt{PivotManager::handleDisturbedPivot} re-zeroes the pivot's phase ambiguity
states in $\stateV$ and the corresponding rows and columns of $\covM$ in place, and
no handover takes place.  In the baseline DD path \texttt{DdPivotManager} applies the
transform per $(\text{cons},\text{code})$ group; a group in which the new pivot has
no existing ambiguity is skipped (logged as ``reset-and-reseed''), and its stale
rows are pruned and re-seeded on the new pivot.

The S-basis constraint is realised by freezing the pivot's ambiguity, ionosphere
(and ionosphere-rate, when enabled) and sat-bias covariances to zero and suppressing
their process noise (\texttt{freezePivotAmbiguityCovariance},
\texttt{freezePivotIonoAndBiasCovariance}).  The frozen ambiguity state is the
condition \texttt{isPivotDisturbed} tests.

\texttt{PivotManager::getCleanHandovers} reports the \emph{clean} handovers of the
epoch: those that were neither forced nor triggered by a disturbed pivot.  A clean
handover is a lossless, invertible frame change, so the RTS smoother
inverse-references the affected states across it; at a forced or disturbed handover
the old pivot's ambiguity is discontinuous and the smoother starts a new segment.
The losslessness stated above holds only for the clean subset.

The two modes use distinct code paths that share the algebra above.  In PPP-RTK
the pivot election, the sticky retention policy and the re-reference live in
\texttt{filter/PivotManager}.  The baseline DD path performs the re-reference in
\texttt{filter/DdPivotManager}, under the \texttt{baseline.enable\_pivot\_handover}
flag (default on).  There the S-transform is applied in place, as a row pass over
the state vector and covariance followed by a column pass over the covariance,
without forming the transform matrix; the result passes through the
positive-semidefinite guard \texttt{CovReReference::enforcePsd} (whose only caller is \texttt{DdPivotManager}),
and the ambiguity keys are renamed onto the new pivot.  With the flag off, every
pivot change, including a forced re-rank, erases the DD ambiguities and re-seeds
them.

\subsection{Between-receiver single difference (BRSD)}\label{sec:brsd}

Baseline modes process one or more rovers $\rx_{\mathrm{rov}}$ against a single
fixed-coordinate reference $\rx_{\mathrm{ref}}$; each rover has its own measurement
model.  The BRSD operator forms
\begin{equation}
\brd{\code^{\sat}_{\rx_{\mathrm{rov}}\rx_{\mathrm{ref}},\sigl}}
\;:=\;
\code^\sat_{\rx_{\mathrm{rov}},\sigl} - \code^\sat_{\rx_{\mathrm{ref}},\sigl},
\end{equation}
and analogously for $\brd{\phaseo}$.  The cancellations of $\brd{\cdot}$ are dual to
those of $\bsd{\cdot}$: BRSD cancels the satellite-side common terms
($\sLight\,\mathrm{d}t^\sat$, $B^\sat_\sigl$, $b^\sat_\sigl$) and retains the
difference of the two stations' receiver clocks and biases.  Once $\brd{}$ is
applied, the satellite biases need no state in the filter.  \texttt{BrsdTransform::apply} carries only
the rover-side line of sight in its output: in baseline mode only the rover position
is estimated, so the reference-station line of sight multiplies no state.

\subsection{Double difference (DD)}\label{sec:dd}

The composition $\dd{\cdot} = \bsd{\big(\brd{\cdot}\big)}$ is
order-independent up to noise correlation: $\bsd{\big(\brd{\cdot}\big)} = \brd{\big(\bsd{\cdot}\big)}$.
Applied to the carrier-phase model, with the common-mode terms cancelled, it gives
\begin{equation}
\dd{\phaseo^{\sat,\piv{g}}_{\rx_{\mathrm{rov}}\rx_{\mathrm{ref}},\sigl}}
\;=\;
\dd{\geomrange^{\sat,\piv{g}}_{\rx_{\mathrm{rov}}\rx_{\mathrm{ref}}}}
+ \dd{T^{\sat,\piv{g}}_{\rx_{\mathrm{rov}}\rx_{\mathrm{ref}}}}
- \mu_\sigl\,\dd{I^{\sat,\piv{g}}_{\rx_{\mathrm{rov}}\rx_{\mathrm{ref}},1}}
+ \lambda_\sigl\,\dd{N^{\sat,\piv{g}}_{\rx_{\mathrm{rov}}\rx_{\mathrm{ref}},\sigl}}
+ \dd{\varepsilon^\phaseo}.
\label{eq:dd-phase}
\end{equation}
\Cref{eq:dd-phase} retains the geometry term (linear in the rover position
perturbation), the differential atmospheric delays and the integer DD ambiguity.
This structure is the basis of carrier-phase RTK.

\paragraph{Short vs.\ long baselines.}
On short baselines (\,$\lesssim \SI{10}{km}$\,) the differential atmospheric terms
in \Cref{eq:dd-phase} are negligible at the millimetre level and are not modelled.
The differential ionosphere $\dd{I}$ grows linearly with separation and reaches
typically \SI{1}{m} at a few hundred kilometres, which exceeds the carrier
wavelength.  \texttt{GnssProcessor} offers two mitigations, both off by default:
\begin{itemize}[leftmargin=*, itemsep=0pt]
\item \texttt{baseline.use\_iono\_free} collapses the per-band DD rows into the
      dual-frequency ionosphere-free combination.  It eliminates the first-order
      ionospheric term at the cost of a non-integer ambiguity (IF phase ambiguities
      stay float and are excluded from AR) and $\sim 3\times$ amplified code noise.
      Satellite pairs with fewer than two shared bands are dropped in this mode.
\item \texttt{baseline.estimate\_differential\_tropo} adds a single
      $\Delta\zwd_{\rx_{\mathrm{rov}}\rx_{\mathrm{ref}}}$ random-walk state
      (\texttt{differential\_tropo\_init\_std\_m}, default \SI{0.1}{m};
      \texttt{differential\_tropo\_process\_noise\_m\_per\_sqrt\_s}, default
      \num{1e-4}).  It absorbs the differential wet delay that survives the BRSD on
      long baselines or baselines with a large height difference.
\end{itemize}

\subsection{S-transform structure}

BSSD and DD both have a matrix form
$\stateV_{\mathrm{diff}} = \bm{S}\,\stateV_{\mathrm{undiff}}$, with
$\bm{S}$ a block design transform.  For BSSD restricted to the ambiguity block of one
constellation $g$ on one signal $\sigl$, with $|\satset{g}|$ tracked satellites,
$\bm{S}$ is a
$(|\satset{g}|-1) \times |\satset{g}|$ matrix of the form
\begin{equation}
\bm{S}_{\mathrm{BSSD}}
\;=\;
\begin{bmatrix}
1 & 0 & \cdots & 0 & -1 & 0 & \cdots & 0 \\
0 & 1 & \cdots & 0 & -1 & 0 & \cdots & 0 \\
\vdots & & \ddots & & \vdots & & \\
0 & 0 & \cdots & 0 & -1 & 0 & \cdots & 1
\end{bmatrix},
\end{equation}
where the column of $-1$'s sits at the index of the pivot $\piv{g}$.  The DD
operator is the composition of the two single-difference operators,
$\bm{S}_{\mathrm{DD}} = \bm{S}_{\mathrm{BSSD}}\,\bm{S}_{\mathrm{BRSD}}$, realised in
code as \texttt{BrsdTransform::apply} followed by
\texttt{BssdTransform::differenceBrsd}
(\texttt{include/measurement/BssdTransform.hpp}).  The composition coincides with the
Kronecker product $\mathbf{D}_{\mathrm{sat}} \otimes \mathbf{D}_{\mathrm{rx}}$ of the
between-satellite and between-receiver differencing matrices acting on the vectorised
observations only on a fully populated satellite\,$\times$\,receiver observation
grid.  The Kronecker form is a model; real data give a non-rectangular grid, and the
code differences whatever pairs exist.  The transformed covariance
$\covM_{\mathrm{diff}} = \bm{S}\,\covM_{\mathrm{undiff}}\,\bm{S}^{\Tr}$ is in general
dense even when $\covM_{\mathrm{undiff}}$ is diagonal, because every differenced row
shares the pivot column.  A DD filter therefore cannot be implemented as independent
per-pair scalar updates: the correlated joint update is what integer ambiguity
recovery through the LAMBDA decorrelation requires.
\textcite{teunissen1985generalised} gives the general theory and a catalogue of
admissible S-transforms; \textcite{deJonge1998processing} treats the rank defect.

\paragraph{Implementation notes.}
\texttt{GnssProcessor} does not materialise the chosen S-system as an explicit
minimum-constraint $\bm{S}$-basis matrix.  It is realised implicitly by the pivot-differencing operator
(\texttt{measurement/BssdTransform}), an S-basis-invariant realisation of the
uncombined parameterisation of \textcite{odijk2003fast}.  A
\texttt{DifferencingMode} strategy in \texttt{measurement/MeasurementBuilder}
dispatches the three regimes (undifferenced, BSSD and DD) to separate builder
classes.

Integer recovery on a DD row requires all four contributing observations to share
the same RINEX signal code.  \texttt{BssdTransform::differenceBrsd} marks a row
\texttt{integer\_capable} only when three conjuncts hold: \texttt{!brsd\_aliased} on
both the $\sat_i$ and the pivot BRSD entries (rover and reference agree on each, i.e.\
no cross-attribute alias), and the rover uses the same code on $\sat_i$ as on the
pivot ($\texttt{rov\_code}_i = \texttt{rov\_code}_{\text{piv}}$).  The third conjunct
is independent of the \texttt{brsd\_aliased} flags: a row can be un-aliased on both
entries and still carry a non-integer DCB/IFB residual if the rover mixes codes
across $\sat_i$ and the pivot.  Rows that fail any conjunct are carried as float and
excluded from ambiguity resolution.
